
\renewcommand{\baselinestretch}{2}

\markright{ \hbox{\footnotesize\rm Statistica Sinica
  }\hfill\\[-13pt]
  \hbox{\footnotesize\rm
  }\hfill }

\markboth{\hfill{\footnotesize\rm AUTHOR(S)} \hfill}
{\hfill {\footnotesize\rm WAVELET ADDITIVE FUNCTIONAL COEFFICIENTS} \hfill}

\renewcommand{\thefootnote}{}
$\ $\par

%%%%%%%%%%%%%%%%%%%%%%%%%%%%%%%%%%%%%%%%%%%%%%%%%%%%%%%%%%%%%%%%%%%%%%%%%%%%%%

\fontsize{12}{14pt plus.8pt minus .6pt}\selectfont \vspace{0.8pc}
\centerline{\large\bf SPARSE WAVELET ESTIMATION OF}
\vspace{2pt}
\centerline{\large\bf ADDITIVE FUNCTIONAL COEFFICIENTS}
\vspace{.4cm}
\centerline{Author(s)}
\vspace{.4cm}
\centerline{\it Affiliation(s)}
\vspace{.55cm} \fontsize{9}{11.5pt plus.8pt minus.6pt}\selectfont

%%%%%%%%%%%%%%%%%%%%%%%%%%%%%%%%%%%%%%%%%%%%%%%%%%%%%%%%%%%%%%%%%%%%%%%%%%%%%%

\begin{quotation}
\noindent {\it Abstract:}
We study a varying coefficient model in which the effect of each linear
covariate is additive in several modulating covariates, and each additive
component is expanded in a periodized wavelet basis. The resulting design is
an approximation space of products between a linear covariate and a wavelet evaluated at a
modulator, and the components are estimated by a single 
block lasso problem in
which the levels of the coefficients are left unpenalized and the
wavelet coefficients of each resolution level are grouped in chunks of order
$\log n$. Three results are
established. The population Gram matrix of the product design has its
smallest eigenvalue bounded away from zero under a conditional second moment
condition on the linear covariates and a bounded joint density for the
modulators, with no independence between the two groups of covariates; the
estimator satisfies a nonasymptotic oracle inequality that needs no cone
condition, because the unpenalized levels are profiled out, and
that holds for any weights of bounded ratio; and, at a
resolution level chosen inside an explicit window, the prediction error and
the componentwise error attain the rate of nonlinear wavelet approximation, with no logarithmic factor when the Besov
integrability index is at least two, which is minimax optimal, under the same Besov condition that governs
linear approximation. The last result rests on   a
bound on the ideal risk over chunks that the Besov condition already implies,
so that adaptivity costs no additional assumption. A threshold on the
norms of the fitted components recovers which modulators act on which
coefficients, and the thresholded estimator keeps the rate of the fit. The
estimator is a single convex
program and is fit by standard
group lasso software.
In a simulation study the thresholded estimator recovers
the structure exactly in every replicate from $n = 1000$ and, where the
components are spatially inhomogeneous, predicts better than penalized
splines and within 8\% of an estimator that knows the structure; on hourly
air pollution data it predicts within 3\% of penalized splines in London and
8\% better in Beijing.
An implementation in R and the code that reproduces every
numerical result of the paper are made available with it.

\vspace{9pt}
\noindent {\it Key words and phrases:}
Additive coefficient model, Besov space,  group
lasso, lasso, oracle
inequality, thresholding, varying coefficient model, wavelets.
\par
\end{quotation}\par

\def\thefigure{\arabic{figure}}
\def\thetable{\arabic{table}}

\renewcommand{\theequation}{\thesection.\arabic{equation}}
% o contador de equacoes do template nao reinicia por secao, o que produzia
% "(3.6)" logo depois de "(2.5)"; o \numberwithin reinicia.
\numberwithin{equation}{section}

\fontsize{12}{14pt plus.8pt minus .6pt}\selectfont

%%%%%%%%%%%%%%%%%%%%%%%%%%%%%%%%%%%%%%%%%%%%%%%%%%%%%%%%%%%%%%%%%%%%%%%%%%%%%%
